\documentclass[tikz, border=12pt]{standalone}
\usepackage{tikz}
\usetikzlibrary{shapes.geometric, arrows.meta, positioning, calc, shadows.blur}
\usepackage{xcolor}

\begin{document}

\begin{tikzpicture}[
    font=\sffamily,
    box/.style={
        rectangle, 
        draw=black, 
        thick, 
        rounded corners=3pt, 
        minimum width=5.5cm, 
        minimum height=0.9cm, 
        align=center, 
        fill=white,
        drop shadow={shadow xshift=0.5mm, shadow yshift=-0.5mm, opacity=0.15}
    },
    excl/.style={
        rectangle, 
        draw=black!60, 
        thick, 
        rounded corners=3pt, 
        minimum width=4.5cm, 
        minimum height=0.9cm, 
        align=center, 
        fill=gray!10,
        drop shadow={shadow xshift=0.5mm, shadow yshift=-0.5mm, opacity=0.10}
    },
    arrow/.style={thick, -{Stealth[length=3mm, width=2.5mm]}},
    label/.style={font=\sffamily\small\itshape, align=center}
]

% ===== LEFT COLUMN: MAIN FLOW =====

% Sampling Design
\node[box, minimum width=6.5cm, fill=cyan!8] (sampling) at (0, 9.5) {%
    \textbf{Sampling Design}\\%
    Two-stage stratified (PPS)\\%
    2011 Population Census\\%
    Stage 1: 675 EA\\%
    Stage 2: 45 households/EA
};

% Study population
\node[box, fill=teal!10] (population) at (0, 7.0) {%
    \textbf{Study population}\\%
    Currently pregnant women aged 15--49\\%
    \textcolor{teal!70!black}{\textbf{n = 1,719}}
};

% Complete cases analysis
\node[box, fill=orange!8] (complete) at (0, 4.5) {%
    \textbf{Complete cases analysis}\\%
    Based on outcome variable
};

% Final sample
\node[box, minimum width=5.5cm, fill=green!10] (final) at (0, 2.0) {%
    \textbf{Final analytical sample}\\%
    \textcolor{green!60!black}{\textbf{n = 4,370}}
};

% ===== RIGHT COLUMN: EXCLUSIONS =====

% Exclusion 1
\node[excl] (excl1) at (6.5, 7.0) {%
    \textbf{Excluded: 28,359}\\%
    Not currently pregnant
};

% Exclusion 2
\node[excl] (excl2) at (6.5, 4.5) {%
    \textbf{Excluded: 25,630}\\%
    Missing/other v626a categories
};

% Exclusion 3
\node[excl, fill=gray!5, dashed] (excl3) at (6.5, 2.0) {%
    \textbf{Excluded: (merged)}
};

% ===== ARROWS: MAIN FLOW =====

% Sampling → Population
\draw[arrow] (sampling) -- (population);

% Population → Complete cases
\draw[arrow] (population) -- (complete);

% Complete cases → Final
\draw[arrow] (complete) -- (final);

% ===== ARROWS: EXCLUSIONS =====

% Population → Excl1
\draw[arrow, color=teal!60!black] (population.east) -- (excl1.west);

% Complete → Excl2
\draw[arrow, color=orange!60!black] (complete.east) -- (excl2.west);

% Final → Excl3
\draw[arrow, dashed, color=gray!60] (final.east) -- (excl3.west);

% ===== FINAL SPLIT =====

% Unmet need
\node[box, minimum width=4.5cm, fill=rose!12] (unmet) at (1.5, -1.5) {%
    \textbf{Unmet need}\\%
    \textcolor{red!70!black}{\textbf{1,791 (41.0\%)}}
};

% No unmet need
\node[box, minimum width=4.5cm, fill=blue!10] (nomet) at (-3.5, -1.5) {%
    \textbf{No unmet need}\\%
    \textcolor{blue!70!black}{\textbf{2,579 (59.0\%)}}
};

% Arrows from Final to Unmet/No unmet
\draw[arrow, color=green!60!black] (final.south) -- ++(0, -0.5) -| (unmet.north);
\draw[arrow, color=green!60!black] (final.south) -- ++(0, -0.5) -| (nomet.north);

% ===== LABELS =====

% Exclusion labels with background
\node[fill=white, inner sep=2pt, font=\sffamily\small\bfseries, text=teal!70!black] at ($(population.east)!0.5!(excl1.west)+(0,0.35)$) {Excluded};
\node[fill=white, inner sep=2pt, font=\sffamily\small\bfseries, text=orange!70!black] at ($(complete.east)!0.5!(excl2.west)+(0,0.35)$) {Excluded};

\end{tikzpicture}

\end{document}